\chapter{Binarne relacije}

Relacije koje su naj\v ce\v s\'ce u upotrebi u ``svakodnevnoj matematici'' kao \v sto su $=$ ili $\le$ su binarne.
Podsetimo se da je \emph{binarna relacija na skupu $A$} svaki podskup $\rho \subseteq A \times A$.
Ako je na skupu $A$ definisana binarna relacija $\rho$, umesto $(a, b) \in \rho$ \v cesto, u skladu sa na\v sim navikama, pi\v semo
$a \mathrel\rho b$.

\begin{EX}\label{ft07.ex.equiv}
  Neka je $m \ge 2$ prirodan broj. Na skupu $\ZZ$ smo definisali binarnu relaciju $\equiv_m$ ovako:
  \begin{center}
    $a \equiv_m b$ ako i samo ako je $a - b$ deljivo sa $m$.
  \end{center}
  Kao \v sto vidimo, i ovde je obi\v caj da pi\v semo $a \equiv_m b$ umesto $(a, b) \in \Boxed{\equiv_m}$.
\end{EX}

\begin{EX}
  Relacija $\le$ na skupu $\NN$ je jedna binarna relacija. Uobi\v cajeno je da pi\v semo $3 \le 7$ umesto $(3, 7) \in \Boxed\le$.
\end{EX}

U nastavku posebnu pa\v znju posve\'cujemo uop\v stenjima relacije jednakosti i relacije poretka na brojevima.
Uop\v stavanjem relacije jednakosti dolazimo do pojma \emph{relacije ekvivalencije} (koja ka\v ze da su dva
objekta ``ujedna\v cena u odnosu na neko apstraktno svojstvo''), odnosno, do pojma \emph{relacije pretka} koji nam
omogu\'cuje da poredimo i objekte koji nisu brojevi.

\clearpage
\section{Relacije ekvivalencije}

Relacija jednakosti $=$ je jedna od osnovnih relacija u matematici. Ona ima slede\'ce osobine:
\begin{itemize}\itemsep -1.5pt
\item $x = x$ za sve $x$,
\item ako je $x = y$ onda je $y = x$, i
\item ako je $x = y$ i $y = z$ onda je $x = z$.
\end{itemize}
\v Ceste su situacije u kojima \v zelimo da ka\v zemo da su neka dva objekta ujedna\v cena u odnosu na neko svojstvo.
Za objekte koji su ujedna\v ceni u odnosu na neko svojstvo ka\v zemo jo\v s i da su \emph{ekvivalentni} s obzirom na
uo\v ceno svojstvo.

Na primer, brojevi 7 i 22 imaju isti ostatak pri deljenju sa~5. Oni su, stoga, ekvivalentni u odnosu na
deljivost sa~5, i to smo u Primeru~\ref{ft07.ex.equiv} ozna\v cavali ovako: $7 \equiv_5 22$.
Primetimo da i relacija $\equiv_5$ ima tri osobine navedene na po\v cetku odeljka:
\begin{itemize}\itemsep -1.5pt
\item $x \equiv_5 x$ za sve $x \in \ZZ$,
\item ako je $x \equiv_5 y$ onda je $y \equiv_5 x$, i
\item ako je $x \equiv_5 y$ i $y \equiv_5 z$ onda je $x \equiv_5 z$.
\end{itemize}
Ove tri osobine karakteri\v su pojam ekvivalentnosti objekata s obzirom na neko svojstvo, a njihovom
apstrakcijom dobijamo pojam relacije ekvivalencije.

\emph{Relacija ekvivalencije na skupu $A$} je svaka binarna relacija $\Boxed\approx$ na skupu $A$
koja ima slede\'ce tri osobine:
\begin{itemize}\itemsep -1.5pt
\item
  (refleksivnost) $a \approx a$ za sve $a \in A$,
\item
  (simetri\v cnost) ako je $a \approx b$ onda je $b \approx a$, za sve $a, b \in A$, i
\item
  (tranzitivnost) ako je $a \approx b$ i $b \approx c$ onda je $a \approx c$, za sve $a, b, c \in A$.
\end{itemize}
Refleksivnost ka\v ze da je svaki objekat ``jednak'' sebi; simetri\v cnost ka\v ze: ako je $a$ ``jednako'' sa $b$, onda je
$b$ ``jednako'' sa $a$; a tranzitivnost ka\v ze: ako je $a$ ``jednako'' sa $b$ i ako je $b$ ``jednako'' sa $c$, onda je
$a$ ``jednako'' sa $c$. Izborom oznake $a \approx b$ \v zelimo da ka\v zemo da $a$ i $b$ \emph{nisu nu\v zno jednaki},
ve\'c da su ujedna\v ceni u odnosu na neko svojstvo.

\begin{EX}\label{ft07.ex.equiv2}
  U Primeru~\ref{ft07.ex.equiv} smo za svaki prirodan broj $m \ge 2$ definisali binarnu relaciju $\equiv_m$
  na skupu $\ZZ$ ovako: $a \equiv_m b$ ako i samo ako $a$ i $b$ imaju isti ostatak pri deljenju sa $m$.
  Lako se dokazuje da je $\equiv_m$ relacija ekvivalencije:
  \begin{itemize}
  \item
    (refleksivnost) za svaki ceo broj $a$ je $a \equiv_m a$ zato \v sto isti brojevi imaju isti ostatak pri deljenju sa~$m$;
  \item
    (simetri\v cnost) ako je $a \equiv_m b$ onda $a$ i $b$ imaju isti ostatak pri deljenju sa $m$; onda $b$ i $a$ imaju
    isti ostatak pri deljenju sa $m$, pa je $b \equiv_m a$;
  \item
    (tranzitivnost) ako je $a \equiv_m b$ i $b \equiv_m c$ onda $a$ i $b$ imaju isti ostatak pri deljenju sa $m$ i
    $b$ i $c$ imaju isti ostatak pri deljenju sa $m$; dakle sva tri broja imaju isti ostatak pri deljenju sa $m$, pa zato
    $a$ i $c$ imaju isti ostatak pri deljenju sa $m$, odnosno $a \equiv_m c$.
  \end{itemize}
\end{EX}

\begin{EX}\label{ft07.ex.01-weight}
  Neka je $n \ge 1$ prirodan broj i neka je $A = \{0,1\}^n$ skup svih 01-nizova du\v zine $n$. Na skupu $A$ defini\v semo
  relaciju $\sim$ ovako: za 01-nizove $u$ i $v$ stavimo $u \sim v$ ako i samo ako nizovi $u$ i $v$ imaju isti broj jedinica
  (broj jedinica u 01-nizu se zove \emph{te\v zina} tog niza; dakle dva niza su ekvivalentni ako imaju istu te\v zinu).
  Na primer, za $n = 8$ je $00001111 \sim 01010110$ i $00110011 \not\sim 11001000$.
  Lako se vidi da je ovo relacija ekvivalencije na $A$.
\end{EX}

\begin{EX}\label{ft07.ex.iskazne-fle}
  Neka je $\Phi$ skup svih iskaznih formula. Na skupu $\Phi$ smo definisali relaciju $\equiv$ ovako:
  $\phi \equiv \psi$ ako i samo ako $\vDash \phi \Leftrightarrow \psi$. Relacija $\equiv$ je relacija
  ekvivalencije na $\Phi$ zbog slede\'ce tri tautologije:
  \begin{itemize}\itemsep -1.5pt
  \item
    (refleksivnost) $\vDash \phi \Leftrightarrow \phi$ za svaku iskaznu formulu $\phi$;
  \item
    (simetri\v cnost) $\vDash (\phi \Leftrightarrow \psi) \Leftrightarrow (\psi \Leftrightarrow \phi)$;
  \item
    (tranzitivnost) $\vDash (\phi \Leftrightarrow \psi) \land (\psi \Leftrightarrow \theta) \Rightarrow (\phi \Leftrightarrow \theta)$.
  \end{itemize}
\end{EX}

Neka je $\Boxed\approx \subseteq A \times A$ relacija ekvivalencije na $A$ i neka je $a \in A$ proizvoljan element.
Skup svih elemenata skupa $A$ koje je relacija $\approx$ ujedna\v cila sa $a$ zovemo \emph{klasa ekvivalencije elementa $a$}
i ozna\v cavamo ovako:
$$
  [a]_\approx = \{ b \in A : a \approx b \}.
$$

\begin{THM}\label{ft07.thm.ekviv}
  Neka je $\Boxed\approx \subseteq A \times A$ relacija ekvivalencije na $A$.
  \begin{enumerate}
  \item
    Za sve $a \in A$ je $[a]_\approx \ne \0$.
  \item
    Neka je $a \in A$. Ako je $x \in [a]_\approx$ onda je $[x]_\approx = [a]_\approx$.
  \item
    Za sve $a, b \in A$ ili je $[a]_\approx \cap [b]_\approx = \0$ ili je $[a]_\approx = [b]_\approx$.
  \item
    Neka je $A = \{a_1, a_2, \ldots, a_n\}$. Tada je $A = [a_1]_\approx \cup [a_2]_\approx \cup \ldots \cup [a_n]_\approx$.
  \end{enumerate}
\end{THM}
\begin{proof}[Dokaz]
  (1)
  Neka je $a \in A$ proizvoljno. Zbog refleksivnosti je $a \approx a$ pa je $a \in [a]_\approx$. Zato
  $[a]_\approx \ne \0$.
  
  \medskip
  
  (2)
  Neka je $a \in A$ i $x \in [a]_\approx$. To zna\v ci da je $x \approx a$. Da bismo pokazali da je
  $[x]_\approx = [a]_\approx$ dokaza\'cemo da je $[x]_\approx \subseteq [a]_\approx$ i da je
  $[x]_\approx \supseteq [a]_\approx$.
  
  $(\subseteq)$
  Neka je $y \in [x]_\approx$. To zna\v ci da je $y \approx x$. Kako $x \approx a$ zaklju\v cujemo da je
  $y \approx a$, pa je zato $y \in [a]_\approx$.

  $(\supseteq)$
  Neka je $y \in [a]_\approx$. To zna\v ci da je $y \approx a$. Kako $x \approx a$ zaklju\v cujemo da je
  $y \approx x$, pa je zato $y \in [x]_\approx$.
  
  \medskip
  
  (3)
  Neka su $a, b \in A$. Mo\v ze se desiti da je $[a]_\approx \cap [b]_\approx = \0$ i u tom slu\v caju je dokaz
  zavr\v sen. Me\dj utim, ako $[a]_\approx \cap [b]_\approx \ne \0$ onda postoji neko $x \in [a]_\approx \cap [b]_\approx$.
  Tada je $x \in [a]_\approx$ i $x \in [b]_\approx$, pa prema (2) sledi da je
  $[x]_\approx = [a]_\approx$ i $[x]_\approx = [b]_\approx$. Dakle, $[a]_\approx = [b]_\approx$.
  
  \medskip
  
  (4)
  Neka je $A = \{a_1, a_2, \ldots, a_n\}$. Da bismo dokazali da je
  $A = [a_1]_\approx \cup [a_2]_\approx \cup \ldots \cup [a_n]_\approx$ dokaza\'cemo da je
  $A \subseteq [a_1]_\approx \cup [a_2]_\approx \cup \ldots \cup [a_n]_\approx$ i
  $A \supseteq [a_1]_\approx \cup [a_2]_\approx \cup \ldots \cup [a_n]_\approx$.
  
  $(\subseteq)$
  Uzmimo proizvoljno $a_i \in A$. Kako smo videli u (1), $a_i \in [a_i]_\approx$, pa je
  $a_i \in [a_1]_\approx \cup [a_2]_\approx \cup \ldots \cup [a_n]_\approx$.
  
  $(\supseteq)$
  Prema definiciji je $[a_i]_\approx \subseteq A$ za sve $i$, pa onda mora biti i cela unija sadr\v zana u $A$:
  $[a_1]_\approx \cup [a_2]_\approx \cup \ldots \cup [a_n]_\approx \subseteq A$.
\end{proof}

\begin{EX}
  Pogledajmo crte\v z relacije $\equiv_3$ na skupu $\{0, 1, 2, 3, 4, 5, 6, 7, 8\}$:
  \begin{center}
    \input ft07-r4.pgf
  \end{center}
  Jasno se vidi da je
  $[0]_{\equiv_3} = [3]_{\equiv_3} = [6]_{\equiv_3} = \{0, 3, 6\}$,
  $[1]_{\equiv_3} = [4]_{\equiv_3} = [7]_{\equiv_3} = \{1, 4, 7\}$, i
  $[2]_{\equiv_3} = [5]_{\equiv_3} = [8]_{\equiv_3} = \{2, 5, 8\}$.
  Drugim re\v cima, skup $\{0, 1, 2, 3, 4, 5, 6, 7, 8\}$ je ``razbijen'' na tri ``klike'', a svaku kliku \v cine elementi
  koje je relacija $\equiv_3$ ujedna\v cila.
\end{EX}

Ovaj primer predstavlja samo jednu instancu op\v steg fenomena koga opisuje pojam koli\v cni\v ckog skupa relacije ekvivalencije.
Neka je $\approx$ relacija ekvivalencije na skupu $A$. \emph{Koli\v cni\v cki skup} ove relacije je skup \v ciji element
su ``klike'' koje nastaju ujedna\v cavanjem elemenata skupa $A$ u odnosu na relaciju~$\approx$ i ozna\v cava se ovako:
$$
  A / \Boxed\approx = \{[a]_\approx : a \in A\}.
$$
\begin{EX}
  U Primeru~\ref{ft07.ex.equiv} smo za svaki prirodan broj $m \ge 2$ definisali binarnu relaciju $\equiv_m$
  na skupu $\ZZ$ ovako: $a \equiv_m b$ ako i samo ako $a$ i $b$ imaju isti ostatak pri deljenju sa $m$.
  Potom smo u Primeru~\ref{ft07.ex.equiv2} pokazali da je $\equiv_m$ relacija ekvivalencije.
  Posmatrajmo sada relaciju $\equiv_5$. Ona ujedna\v cava elemente skupa $\ZZ$ ako imaju isti ostatak pri deljenju
  sa~5. Na primer,
  \begin{align*}
    [0]_{\equiv_5} &= \{\ldots, -10, -5, 0, 5, 10, \ldots\},\\
    [1]_{\equiv_5} &= \{\ldots, -9, -4, 1, 6, 11, \ldots\},\\
    [2]_{\equiv_5} &= \{\ldots, -8, -3, 2, 7, 12, \ldots\},\\
    [3]_{\equiv_5} &= \{\ldots, -7, -2, 3, 8, 13, \ldots\},\\
    [4]_{\equiv_5} &= \{\ldots, -6, -1, 4, 9, 14, \ldots\}.
  \end{align*}
  Po\v sto pri deljenju sa~5 mo\v zemo dobiti samo pet razli\v citih ostataka (0, 1, 2, 3, 4), u ovih pet skupova
  su svrstani svi celi brojevi pa zaklju\v cujemo da je:
  $$
    \ZZ / \Boxed{\equiv_5} = \{[0]_{\equiv_5}, [1]_{\equiv_5}, [2]_{\equiv_5}, [3]_{\equiv_5}, [4]_{\equiv_5}\}.
  $$
  Uop\v steno, za $m \ge 2$ je
  $$
    \ZZ / \Boxed{\equiv_m} = \{[0]_{\equiv_m}, [1]_{\equiv_m}, [2]_{\equiv_m}, \ldots, [m-1]_{\equiv_m}\}
  $$
  jer pri deljenju sa $m$ mo\v zemo dobiti samo $m$ razli\v citih ostataka (0, 1, 2, \ldots, $m-1$).
\end{EX}

\begin{EX}
  U Primeru~\ref{ft07.ex.01-weight} smo na skupu $A = \{0,1\}^n$ skup svih 01-nizova du\v zine $n$, $n \ge 1$,
  definisali relaciju $\sim$ ovako: $u \sim v$ ako i samo ako nizovi $u$ i $v$ imaju isti broj jedinica.
  Neka je $n = 8$. Primetimo, prvo, da je $00000001 \sim 00000010 \sim 00000100 \sim \ldots \sim 10000000$,
  potom $00000011 \sim 00000101 \sim \ldots \sim 11000000$, kao i da je $00000000$ jedini 01-niz $u \in A^8$
  sa osobinom $u \sim 00000000$. Zato je:
  \begin{align*}
    [00000000]_\sim &= \{00000000\}\\
    [00000001]_\sim &= \{ u \in A^8 : u \text{ ima ta\v cno jednu jedinicu} \},\\
    [00000011]_\sim &= \{ u \in A^8 : u \text{ ima ta\v cno dve jedinice} \},\\
                    &\vdots\\
    [01111111]_\sim &= \{ u \in A^8 : u \text{ ima ta\v cno sedam jedinica} \},\\
    [11111111]_\sim &= \{ 11111111 \}.
  \end{align*}
  Ovim su obuhva\'ceni svi 01-nizovi iz $A^8$, pa zaklju\v cujemo da je
  $$
    A^8 / \Boxed\sim = \{[00000000]_\sim, [00000001]_\sim, [00000011]_\sim, \ldots, [11111111]_\sim\}.
  $$
  Ovaj primer se mo\v ze uop\v stiti na proizvoljno $n \ge 1$. Lako vidimo da je
  \begin{align*}
    [000\ldots00]_\sim &= \{000\ldots00\}\\
    [000\ldots01]_\sim &= \{ u \in A^n : u \text{ ima ta\v cno jednu jedinicu} \},\\
                    &\vdots\\
    [011\ldots11]_\sim &= \{ u \in A^n : u \text{ ima ta\v cno $n-1$ jedinica} \},\\
    [111\ldots11]_\sim &= \{ 111\ldots11 \}.
  \end{align*}
  Skup $A^n / \Boxed\sim$ ima $n + 1$ elemenata i to su ta\v cno ovih $n + 1$ klasa ekvivalencije.
\end{EX}

\begin{EX}
  U Primeru~\ref{ft07.ex.iskazne-fle} smo na skupu $\Phi$ svih iskaznih formula uo\v cili relaciju
  ekvivalencije $\equiv$ ovako: $\phi \equiv \psi$ ako i samo ako $\vDash \phi \Leftrightarrow \psi$.
  Sada se lako vidi da je \hbox{$[p \lor \lnot p]_\equiv$} skup svih tautologija, dok je
  \hbox{$[p \land \lnot p]_\equiv$} skup svih kontradikcija (\emph{kontradikcija} je iskazna formula koja je neta\v cna
  za sve valuacije).
\end{EX}

Videli smo u primerima da relacija ekvivalencije ``razbija'' skup na kome je definisana na ``klike'', gde se u istoj
kliki nalaze elementi koje je ta relacija ujedna\v cila. Sada \'cemo formalizovati ovo zapa\v znje.

  Neka je $A$ skup i neka su $B_1$, $B_2$, \ldots, $B_k$ neki podskupovi skupa $A$.
  Skup $\{B_1, B_2, \ldots, B_k\}$ je \emph{particija} ili \emph{razbijanje} skupa $A$
  ako va\v zi slede\'ce:

\bigskip

\noindent
\begin{minipage}{3in}
  \begin{itemize}\itemsep -1.5pt
  \item $B_i \ne \0$ za sve $i$;
  \item $B_i \cap B_j = \0$ kadgod je $i \ne j$; i
  \item $B_1 \cup B_2 \cup \ldots \cup B_k = A$.
  \end{itemize}

  Skupovi $B_1$, $B_2$, \ldots, $B_k$ se zovu \emph{blokovi} particije.

  \hspace*{1.5em}Dakle, blokovi particije su neprazni podskupovi skupa $A$ sa osobinom da su svaka dva razli\v cita bloka disjunktni
  i da svaki element skupa $A$ pripada nekom bloku (pa zbog druge osobine dobijamo da svaki element skupa $A$ pripada
  ta\v cno jednom bloku).
\end{minipage}
\hfill
\begin{minipage}{2.5in}
  \begin{center}
    \input ft07-r5.pgf
  \end{center}
\end{minipage}

\begin{THM}
  Neka je $A$ kona\v can skup.
  \begin{enumerate}
  \item Ako je $\approx$ relacija ekvivalencije na $A$ onda je $A / \Boxed\approx$ particija skupa $A$.
  \item Obrnuto, za svaku particiju $\{B_1, B_2, \ldots, B_k\}$ skupa $A$ postoji relacija ekvivalencije $\approx$ takva da je
  $A / \Boxed\approx = \{B_1, B_2, \ldots, B_k\}$.
  \end{enumerate}
\end{THM}
\begin{proof}[Dokaz]
  (1)
  Neka je $\approx$ relacija ekvivalencije na $A$. Neka se u skupu $A / \Boxed\approx$ nalazi $k$ po parovima
  razli\v citih klasa. Iz svake klase odaberimo po jednog predstavnika. Neka su to $a_1, a_2, \ldots, a_k \in A$.
  Tada je
  $$
    A / \Boxed\approx = \{[a_1]_\approx, [a_2]_\approx, \ldots, [a_k]_\approx\}.
  $$
  Kako smo videli u Teoremi~\ref{ft07.thm.ekviv}, svi $[a_i]_\approx$ su neprazni, svake dve razli\v cite klase su disjunktne
  i unija svih je~$A$. Dakle, $A / \Boxed\approx$ je particija skupa~$A$.
  
  \medskip
  
  (2)
  Neka je $\{B_1, B_2, \ldots, B_k\}$ particija skupa $A$ i neka je $\approx$ binarna relacija na $A$ definisana ovako:
  $$
    a \approx b \text{\ \ ako i samo ako $a$ i $b$ pripadaju istom bloku particije.}
  $$
  Doka\v zimo da je to relacija ekvivalencije.
  \begin{itemize}
  \item
    (refleksivnost)
    Neka je $a \in A$ proizvoljno. Tada se $a$ nalazi u nekom bloku $B_i$, pa se $a$ i $a$ nalaze u istom bloku.
    Zato je $a \approx a$.
  \item
    (simetri\v cnost)
    Neka je $a \approx b$. Tada se $a$ i $b$ nalaze u istom bloku, pa se onda i $b$ i $a$ nalaze u istom bloku.
    Zato je $b \approx a$.
  \item
    (tranzitivnost)
    Neka je $a \approx b$ i $b \approx c$. Tada se $a$ i $b$ nalaze u istom bloku, i $b$ i $c$ se nalaze u istom bloku.
    Recimo da je $a, b \in B_i$ i $b, c \in B_j$. Zato \v sto su razli\v citi blokovi particije disjunktni, zaklju\v cujemo
    da je $B_i = B_j$ jer $b \in B_i \cap B_j$. Dakle, $a$ i $c$ se tako\dj e nalaze u istom bloku, pa je zato $a \approx c$.
  \end{itemize}
  Treba jo\v s pokazati da je $A / \Boxed\approx = \{B_1, B_2, \ldots, B_k\}$, ali to sledi direktno iz konstrukcije relacije~$\approx$.
\end{proof}

\section{Relacije poretka}

Relacije poretka nam omogu\'cuju da uredimo neki skup po nekom svojstvu. Recimo, brojeve mo\v zemo da
pore\dj amo po veli\v cini, ili neki skup re\v ci mo\v zemo da pore\dj amo po \emph{leksikografskom poretku}, recimo ovako:
\begin{quote}
  Avala,
  Besna kobila,
  Cer,
  Dukat,
  Fruška gora,
  Golija,
  Homoljske planine,
  Jelova gora,
  Kopaonik,
  Lepa gora,
  Maljen,
  Ozren,
  Povlen,
  Rtanj,
  Stara planina,
  Tara,
  Vršačka brda,
  Zlatibor
\end{quote}
Do apstraktnog pojma relacije poretka dolazimo analizom apstraktnih svojstava dva prirodna poretka:
poretka na brojevima ($\le$) i relacije podskupa ($\subseteq$).

Poredak na brojevima $\le$ je jedna od osnovnih relacija u matematici. Ona ima slede\'ce osobine:
\begin{itemize}\itemsep -1.5pt
\item $x \le x$ za svaki broj $x$,
\item ako je $x \le y$ i $y \le x$ onda je $x = y$, i
\item ako je $x \le y$ i $y \le z$ onda je $x \le z$.
\end{itemize}
Interesantno je da i relacija $\subseteq$ na skupovima ima iste ove tri osobine:
\begin{itemize}\itemsep -1.5pt
\item $X \subseteq X$ za svaki skup $X$,
\item ako je $X \subseteq Y$ i $Y \subseteq X$ onda je $X = Y$, i
\item ako je $X \subseteq Y$ i $Y \subseteq Z$ onda je $X \subseteq Z$.
\end{itemize}
Ove tri osobine karakteri\v su poredak objekata, a njihovom
apstrakcijom dobijamo pojam apstraktne relacije poretka.

\emph{Relacija poretka na skupu $A$} je svaka binarna relacija $\Boxed\sqsubseteq$ na skupu $A$
koja ima slede\'ce tri osobine:
\begin{itemize}\itemsep -1.5pt
\item
  (refleksivnost) $a \sqsubseteq a$ za sve $a \in A$,
\item
  (antisimetri\v cnost) ako je $a \sqsubseteq b$ i $b \sqsubseteq a$ onda je $a = b$, za sve $a, b \in A$, i
\item
  (tranzitivnost) ako je $a \sqsubseteq b$ i $b \sqsubseteq c$ onda je $a \sqsubseteq c$, za sve $a, b, c \in A$.
\end{itemize}
Refleksivnost ka\v ze da je svaki objekat ``manji ili jednak'' od sebe;
antisimetri\v cnost ka\v ze: ako je $a$ ``manje ili jednako'' sa $b$ i ako je $b$ ``manje ili jednako'' sa $a$ onda je
$a = b$; a tranzitivnost ka\v ze: ako je $a$ ``manje ili jednako'' sa $b$ i ako je $b$ ``manje ili jednako'' sa $c$, onda je
$a$ ``manje ili jednako'' sa $c$.

Neka je $\le$ uobi\v cajeni poredak celih brojeva. Ova relacija poretka osim tri navedene ima jo\v s jednu
interesantnu osobinu: \emph{svaka dva cela broja su uporediva}.
Naime, za svaka dva broja $a, b \in \ZZ$ je $a \le b$ ili je $b \le a$.
S druge strane, poredak $\subseteq$ na skupovima nema tu osobinu: postoje skupovi koji su \emph{neuporedivi}.
Na primer $\{1, 2\} \not\subseteq \{2, 3\}$ i $\{2, 3\} \not\subseteq \{1, 2\}$.

Relacija poretka na skupu $A$ je \emph{linearni (ili totalni) poredak} ako va\v zi slede\'ce
\begin{itemize}\itemsep -1.5pt
\item
  za svaka dva elementa $a, b \in A$ je $a \sqsubseteq b$ ili je $b \sqsubseteq a$.
\end{itemize}

\begin{EX}\label{ft07.ex.deljivost}
  Neka je $\mid$ relacija deljivosti na skupu $\NN$. Na primer, $3 \mid 12$ i $5 \nmid 12$. Tada je lako dokazuje da je
  $\mid$ relacija poretka na $\NN$. Refleksivnost i tranzitivnost su o\v cigledni. Doka\v zimo antisimetri\v cnost.
  Neka je $a \mid b$ i $b \mid a$ za neke $a, b \in \NN$. Iz $a \mid b$ sledi $a \le b$. Isto tako, iz $b \mid a$ sledi
  $b \le a$. Dakle, $a \le b$ i $b \le a$, pa je $a = b$. Ovo nije linearni poredak jer postoje neuporedivi elementi,
  recimo, $2 \nmid 3$ i $3 \nmid 2$.
\end{EX}

\begin{EX}
  Neka je na skupu $\NN \times \NN$ definisana relacija $\sqsubseteq$ ovako:
  $$
    (a, b) \sqsubseteq (c, d) \text{\ \ ako i samo ako\ \ } a \le c \land b \le d.
  $$
  Tada se lako vidi da je $\sqsubseteq$ relacija poretka. Ovo nije linearni poredak jer postoje neuporedivi elementi,
  recimo, $(1, 2) \not\sqsubseteq (2, 1)$ i $(2, 1) \not\sqsubseteq (1, 2)$.
\end{EX}

Prethodni primer je specijalni slu\v caj slede\'ceg rezultata koji pokazuje kako mo\v ze da se konstrui\v se
relacija poretka na proizvodu.

\begin{THM}
  Neka je $\sqsubseteq_1$ relacija poretka na skupu $A_1$, a $\sqsubseteq_2$ relacija poretka na skupu $A_2$.
  Na skupu $A_1 \times A_2$ defini\v semo binarnu relaciju $\preccurlyeq$ ovako:
  $$
    (a_1, a_2) \preccurlyeq (b_1, b_2) \text{\ \ ako i samo ako je\ \ } a_1 \mathrel{\sqsubseteq_1} b_1 \land a_2 \mathrel{\sqsubseteq_2} b_2.
  $$
  Tada je $\preccurlyeq$ relacija poretka na $A_1 \times A_2$.
\end{THM}
\begin{proof}[Dokaz]
  Doka\v zimo da je relacija $\preccurlyeq$ refleksivna, antisimetri\v cna i tranzitivna. Primetimo, prvo, da je
  $$
    (a_1, a_2) \preccurlyeq (a_1, a_2) \Leftrightarrow a_1 \mathrel{\sqsubseteq_1} a_1 \land a_2 \mathrel{\sqsubseteq_2} a_2.
  $$
  Kako je iskaz na desnoj strani ta\v can (zato \v sto su relacije $\sqsubseteq_1$ i $\sqsubseteq_2$ refleksivne)
  pokazali smo da je relacija $\preccurlyeq$ refleksivna. Da bismo pokazali da je relacija $\preccurlyeq$ antisimetri\v cna
  pretpostavimo da je $(a_1, a_2) \preccurlyeq (b_1, b_2)$ i $(b_1, b_2) \preccurlyeq (a_1, a_2)$. Prema definiciji je:
  \begin{align*}
    (a_1, a_2) &\preccurlyeq (b_1, b_2) \land (b_1, b_2) \preccurlyeq (a_1, a_2)\\
               &\Leftrightarrow a_1 \mathrel{\sqsubseteq_1} b_1 \land a_2 \mathrel{\sqsubseteq_2} b_2 \land b_1 \mathrel{\sqsubseteq_1} a_1 \land b_2 \mathrel{\sqsubseteq_2} a_2\\
               &\Leftrightarrow (a_1 \mathrel{\sqsubseteq_1} b_1 \land b_1 \mathrel{\sqsubseteq_1} a_1) \land (a_2 \mathrel{\sqsubseteq_2} b_2 \land b_2 \mathrel{\sqsubseteq_2} a_2)\\
               &\Leftrightarrow a_1 = b_1 \land a_2 = b_2\\
               &\Leftrightarrow (a_1, a_2) \preccurlyeq (b_1, b_2),
  \end{align*}
  zato \v sto su relacije $\sqsubseteq_1$ i $\sqsubseteq_2$ antisimetri\v cne. Kona\v cno da bismo pokazali
  da je relacija $\preccurlyeq$ tranzitivna pretpostavimo da je $(a_1, a_2) \preccurlyeq (b_1, b_2)$ i $(b_1, b_2) \preccurlyeq (c_1, c_2)$.
  Tada je:
  \begin{align*}
    (a_1, a_2) &\preccurlyeq (b_1, b_2) \land (b_1, b_2) \preccurlyeq (c_1, c_2)\\
               &\Leftrightarrow a_1 \mathrel{\sqsubseteq_1} b_1 \land a_2 \mathrel{\sqsubseteq_2} b_2 \land b_1 \mathrel{\sqsubseteq_1} c_1 \land b_2 \mathrel{\sqsubseteq_2} c_2\\
               &\Leftrightarrow (a_1 \mathrel{\sqsubseteq_1} b_1 \land b_1 \mathrel{\sqsubseteq_1} c_1) \land (a_2 \mathrel{\sqsubseteq_2} b_2 \land b_2 \mathrel{\sqsubseteq_2} c_2)\\
               &\Leftrightarrow a_1 \mathrel{\sqsubseteq_1} c_1 \land a_2 \mathrel{\sqsubseteq_2} c_2\\
               &\Leftrightarrow (a_1, a_2) \preccurlyeq (c_1, c_2),
  \end{align*}
  zato \v sto su relacije $\sqsubseteq_1$ i $\sqsubseteq_2$ tranzitivne.
\end{proof}

\begin{EX}
  Defini\v simo na skupu $\NN \times \NN$ relaciju $\preccurlyeq$ ovako:
  $$
    (a, b) \preccurlyeq (c, d) \text{\ \ ako i samo ako\ \ } a < c \lor (a = c \land b \le d).
  $$
  Lako se pokazuje da je ovo relacija poretka. Doka\v zimo da se radi o linearnom poretku. Neka su
  $(a, b), (c, d) \in \NN \times \NN$ proizvoljni. Ako je $a < c$ onda je
  $(a, b) \preccurlyeq (c, d)$; ako je $c \le a$ onda je $(c, d) \preccurlyeq (a, b)$; ako je $a = c$
  onda odnos brojeva $b$ i $d$ odre\dj uje odnos parova $(a, b)$ i $(c, d)$: ako je $b \le d$ onda je
  $(a, b) \preccurlyeq (c, d)$, a ako je $d \le b$ onda je $(c, d) \preccurlyeq (a, b)$. Time smo pokazali da su
  svaka dva elementa skupa $\NN \times \NN$ uporediva.
\end{EX}

Prethodni primer je specijalni slu\v caj slede\'ceg rezultata koji pokazuje kako mo\v ze da se konstrui\v se
relacija poretka na proizvodu.

\begin{THM}\label{ft07.thm.lex-2}
  Neka je $\sqsubseteq_1$ linearni poredak na skupu $A_1$, a $\sqsubseteq_2$ linearni poredak na skupu $A_2$.
  Na skupu $A_1 \times A_2$ defini\v semo binarnu relaciju $\preccurlyeq$ ovako:
  $$
    (a_1, a_2) \preccurlyeq (b_1, b_2) \text{\ \ ako i samo ako je\ \ }
    (a_1 \ne b_1 \land a_1 \mathrel{\sqsubseteq_1} b_1) \lor (a_1 = b_1 \land a_2 \mathrel{\sqsubseteq_2} b_2).
  $$
  Tada je $\preccurlyeq$ linearni poredak na $A_1 \times A_2$.
\end{THM}
\begin{proof}[Dokaz]
  Doka\v zimo, prvo, da je $\preccurlyeq$ relacija poretka na $A_1 \times A_2$. Refleksivnost se pokazuje lako:
  $(a_1, a_2) \preccurlyeq (a_1, a_2)$ je ta\v cno zato \v sto je drugi disjunkt u iskazu sa desne strane ta\v can
  ($a_1 = a_1 \land a_2 \mathrel{\sqsubseteq_2} a_2$).
  
  \medskip
  
  Doka\v zimo sada antisimetri\v cnot. Neka je $(a_1, a_2) \preccurlyeq (b_1, b_2)$ i $(b_1, b_2) \preccurlyeq (a_1, a_2)$.
  Kao prvi korak dokaza\'cemo da je $a_1 = b_1$. Pretpostavimo da to nije ta\v cno. Tada je $a_1 \ne b_1$, pa na osnovu
  definicije dobijamo
  \begin{align*}
    (a_1, a_2) \preccurlyeq (b_1, b_2) &\Rightarrow a_1 \ne b_1 \land a_1 \mathrel{\sqsubseteq_1} b_1 \Rightarrow a_1 \mathrel{\sqsubseteq_1} b_1\\
    (b_1, b_2) \preccurlyeq (a_1, a_2) &\Rightarrow b_1 \ne a_1 \land b_1 \mathrel{\sqsubseteq_1} a_1 \Rightarrow b_1 \mathrel{\sqsubseteq_1} a_1,
  \end{align*}
  (zato \v sto je drugi disjunkt u iskazu sa desne strane neta\v can u oba slu\v caja). Zbog antisimetri\v cnosti relacije $\sqsubseteq_1$
  zaklju\v cujemo da je $a_1 \ne b_1$, \v sto je u suprotnosti sa pretpostavkom $a_1 = b_1$.
  
  Dakle, mora biti $a_1 = b_1$. Sada se lako pokazuje da je $a_2 = b_2$:
  \begin{align*}
    (a_1, a_2) \preccurlyeq (b_1, b_2) &\Rightarrow a_1 = b_1 \land a_2 \mathrel{\sqsubseteq_2} b_2 \Rightarrow a_2 \mathrel{\sqsubseteq_2} b_2\\
    (b_1, b_2) \preccurlyeq (a_1, a_2) &\Rightarrow b_1 = a_1 \land b_2 \mathrel{\sqsubseteq_2} a_2 \Rightarrow b_2 \mathrel{\sqsubseteq_2} a_2,
  \end{align*}
  pa je $a_2 = b_2$. Dakle, $(a_1, b_1) = (a_2, b_2)$.
  
  \medskip
  
  Da bismo pokazali tranzitivnost pretpostavimo da je $(a_1, a_2) \preccurlyeq (b_1, b_2)$ i $(b_1, b_2) \preccurlyeq (c_1, c_2)$.
  Razmotri\'cemo \v cetiri slu\v caja.
  
  \medskip
  
  Slu\v caj 1: $a_1 = b_1$ i $b_1 = c_1$. Tada je $a_2 \mathrel{\sqsubseteq_2} b_2$ i $b_2 \mathrel{\sqsubseteq_2} c_2$.
  Dakle, $a_1 = c_1$ i $a_2 \mathrel{\sqsubseteq_2} c_2$, pa je $(a_1, a_2) \preccurlyeq (c_1, c_2)$.
  
  \medskip

  Slu\v caj 2: $a_1 = b_1$ i $b_1 \ne c_1$. Tada je $b_1 \mathrel{\sqsubseteq_1} c_1$.
  Dakle, $a_1 \ne c_1$ i $a_1 \mathrel{\sqsubseteq_1} c_1$ (jer je $a_1 = b_1$), pa je $(a_1, a_2) \preccurlyeq (c_1, c_2)$.

  \medskip

  Slu\v caj 3: $a_1 \ne b_1$ i $b_1 = c_1$. Tada je $a_1 \mathrel{\sqsubseteq_1} b_1$.
  Dakle, $a_1 \ne c_1$ i $a_1 \mathrel{\sqsubseteq_1} c_1$ (jer je $b_1 = c_1$), pa je $(a_1, a_2) \preccurlyeq (c_1, c_2)$.

  \medskip

  Slu\v caj 4: $a_1 \ne b_1$ i $b_1 \ne c_1$. Tada je $a_1 \mathrel{\sqsubseteq_1} b_1$ i $b_1 \mathrel{\sqsubseteq_1} c_1$.
  Tada je $a_1 \mathrel{\sqsubseteq_1} c_1$. Lako se pokazuje da je $a_1 \ne c_1$
  (jer u suprotnom iz $a_1 = c_1$ sledi $c_1 \le a_1$, pa $a_1 \mathrel{\sqsubseteq_1} b_1$ dalje implicira $c_1 \mathrel{\sqsubseteq_1} b_1$;
  to sa zaklju\v ckom $b_1 \mathrel{\sqsubseteq_1} c_1$ daje $b_1 = c_1$, \v sto je suprotno pretpostavci $b_1 \ne c_1$).
  Dakle, $a_1 \ne c_1$ i $a_1 \mathrel{\sqsubseteq_1} c_1$, pa je $(a_1, a_2) \preccurlyeq (c_1, c_2)$.
  
  \medskip
  
  Kona\v cno, doka\v zimo da je $\preccurlyeq$ linearni poredak, odnosno, da su svaka dva elementa skupa $A_1 \times A_2$
  uporediva. Neka su $(a_1, a_2), (b_1, b_2) \in A_1 \times A_2$ proizvoljni. Razmotri\'cemo dva slu\v caja.
  
  \medskip
  
  Slu\v caj 1: $a_1 = b_1$. Zato \v sto je $\sqsubseteq_2$ linearni poredak mora bili $a_2 \mathrel{\sqsubseteq_2} b_2$
  ili $b_2 \mathrel{\sqsubseteq_2} a_2$, pa je, u zavisnosti od toga koja od ove dve relacije je ta\v cna,
  $(a_1, a_2) \preccurlyeq (b_1, b_2)$ ili $(b_1, b_2) \preccurlyeq (a_1, a_2)$.
  
  \medskip
  
  Slu\v caj 2: $a_1 \ne b_1$. Zato \v sto je $\sqsubseteq_1$ linearni poredak mora bili $a_1 \mathrel{\sqsubseteq_1} b_1$
  ili $b_1 \mathrel{\sqsubseteq_1} a_1$, pa je, u zavisnosti od toga koja od ove dve relacije je ta\v cna,
  $(a_1, a_2) \preccurlyeq (b_1, b_2)$ ili $(b_1, b_2) \preccurlyeq (a_1, a_2)$.
  
  \medskip
  
  Ovim je tvr\dj enje pokazano.
\end{proof}



\begin{EX}
  Da se podsetimo, za kona\v can skup $A$ jo\v s ka\v zemo i da je \emph{alfabet}.
  Za njegove elemente tada ka\v zemo da su \emph{slova}.
  U Primeru~\ref{ft06.ex.A-star} smo videli da je re\v c nad alfabetom $A$ proizvoljan kona\v can niz slova tog alfabeta,
  uklju\v cuju\'ci i praznu re\v c koju smo ozna\v cili sa~$\epsilon$.
  Skup svih re\v ci nad alfabetom $A$ je, dakle, skup svih kona\v cnih nizova elemenata skupa~$A$:
  $$
    A^* = \{\epsilon\} \cup A^1 \cup A^2 \cup \ldots \cup A^n \cup \ldots.
  $$
  Re\v c $u$ je \emph{prefiks} re\v ci $v$ ako se re\v c $u$ nalazi na samom po\v cetku re\v ci $v$, iza \v cega u re\v ci
  $v$ mo\v ze, a ne mora da se nalazi jo\v s slova. Relacija ``\ldots\ je prefiks od \ldots'' je relacija poretka
  na $A^*$ koja ne mora biti linearni poredak. Na primer, za $A = \{a, b, c\}$ imamo slede\'ce:
  \begin{itemize}\itemsep -1.5pt
  \item
    $ab$ je prefiks od $abc$, a nije prefiks od $aab$;
  \item
    prazna re\v c je prefkis svake re\v ci;
  \item
    za svaku re\v c $u$ je $u$ prefiks od $u$ (refleksivnost);
  \item
    ako je $u$ prefiks od $v$ i ako je $v$ prefiks od $u$ onda je $u = v$ (antisimetri\v cnost);
  \item
    ako je $u$ prefiks od $v$ i ako je $v$ prefiks od $w$ onda je $u$ prefiks od $w$ (tranzitivnost);
  \item
    niti je $ab$ prefiks od $ba$, niti je $ba$ prefiks od $ab$ (zato ovo nije linearni poredak).
  \end{itemize}
\end{EX}

\begin{EX}
  Relacija poretka koju smo konstruisali u Teoremi~\ref{ft07.thm.lex-2} se zove \emph{leksikografski poredak} jer se ista ideja
  koristi za ure\dj ivanje re\v ci u leksikonu. Cilj ovog primera je da uop\v stimo navedenu konstrukciju na skup re\v ci nad
  nekim alfabetom.
  
  Neka je $A$ alfabet i neka je $\sqsubseteq$ proizvoljan linearni poredak na $A$. Na skupu $A^*$
  defini\v semo linearni poredak $\preccurlyeq$ ovako. Neka su $u = a_1 a_2 \ldots a_m$ i $v = b_1 b_2 \ldots b_n$ proizvoljne
  re\v ci iz~$A^*$.
  \begin{itemize}
  \item
    Ako je $u$ prefiks od $v$ onda je $u \preccurlyeq v$; ako je $v$ prefiks od $u$ onda je $v \preccurlyeq u$.
  \item
    Ako nijedna od ove dve re\v ci nije prefiks one druge uo\v cimo najmanje $i$ takvo da je $a_i \ne b_i$,
    (to zna\v ci da je $a_1 = b_1$, $a_2 = b_2$, \ldots, $a_{i-1} = b_{i-1}$), pa za $a_i < b_i$ stavimo
    $u \preccurlyeq v$, dok za $b_i < a_i$ stavimo $v \preccurlyeq u$.
  \end{itemize}
  Tada je $\preccurlyeq$ linearni poredak na~$A^*$.
  Na primer, za $A = \{a, b, c, d\}$ imamo slede\'ce: $aaaaa \preccurlyeq b \preccurlyeq bcac \preccurlyeq bcd$.
\end{EX}

Relacije poretka kao skupovi ure\dj enih parova obi\v cno imaju mnogo elemenata pa grafi\v cko predstavljanje koje smo opisali
u Primerima~\ref{ft07.ex.01}, \ref{ft07.ex.02} i~\ref{ft07.ex.03} nije pogodno. Zato se relacije poretka predstavljaju
grafi\v cki koriste\'ci Haseove dijagrame. Haseov dijagram predstavlja relaciju poretka tako \v sto na dijagram unosi minimalnu
koli\v cinu informacija, uz dogovor da se manji element crtaju ispod onih koji su ve\'ci od njih. Ne\'cemo ovde ulaziti
u formalnu raspravu o tome kako se ta\v cno defini\v se Haseov dijagram, ve\'c cemo na nekoliko primera pokazati kako se on
konstrui\v se za datu relaciju poretka.

\begin{EX}\label{ft07.ex.deljivost-1-10}
  Relacija deljivosti $\mid$ je relacija poretka na skupu $\{1, 2, 3, 4, 5, 6, 7, 8, 9, 10\}$
  (Primer~\ref{ft07.ex.deljivost}). Pokaza\'cemo sada kako se crta Haseov dijagram koji joj odgovara.
\end{EX}

\bigskip

\noindent
\begin{minipage}{3in}
  Crtanje kre\'cemo od ``najmanjeg'' elementa, \v sto je u ovom slu\v caju~1.
\end{minipage}
\hfill
\begin{minipage}{2.5in}
  \begin{center}
    \input ft07-r6a.pgf
  \end{center}
\end{minipage}

\bigskip

\noindent
\begin{minipage}{3in}
  Slede\'ci na redu je element~2. Kako je 2 ve\'ce od 1, nacrta\'cemo~2 iznad~1 i spojiti ih linijom.
\end{minipage}
\hfill
\begin{minipage}{2.5in}
  \begin{center}
    \input ft07-r6b.pgf
  \end{center}
\end{minipage}

\bigskip

\noindent
\begin{minipage}{3in}
  Broj 3 je deljiv sa 1, ali nije deljiv sa 2. Zato \'cemo broj 3 nacrtati iznad 1 i spojiti ga linijom sa 1, a ne\'cemo
  ga spojiti linijom sa 2.
\end{minipage}
\hfill
\begin{minipage}{2.5in}
  \begin{center}
    \input ft07-r6c.pgf
  \end{center}
\end{minipage}

\bigskip

\noindent
\begin{minipage}{3in}
  Broj 4 je deljiv brojem 2, ali nije deljiv brojem 3. Zato \'cemo broj 4 nacrtati iznad broja 2 i spojiti ga linijom sa 2, a ne\'cemo
  ga spojiti linijom sa 3.
\end{minipage}
\hfill
\begin{minipage}{2.5in}
  \begin{center}
    \input ft07-r6d.pgf
  \end{center}
\end{minipage}

\medskip

Primetimo da je broj 4 deljiv i brojem 1. Me\dj utim, ne\'cemo broj 4 spojiti linijom sa brojem 1 da bismo smanjili broj linija na
dijagramu. Zato se kod Haseovih dijagrama uvodi slede\'ci dogovor: element $x$ je manji od elementa $y$ ako je $x$ ispod $y$ i postoji
put koji spaja $x$ sa $y$ u kome linije idu stalno navi\v se. Ako takav put ne postoji, onda su $x$ i $y$ neuporedivi.
To je slu\v caj sa brojevima 3 i 4 na poslednjem dijagramu: zato \v sto su 3 i 4 neuporedivi ne postoji put od 3 do 4 koji ide
samo navi\v se.

\bigskip

\noindent
\begin{minipage}{3in}
  Od svih navedenih brojeva broj 5 je deljiv samo brojem 1. Zato \'cemo broj 5 nacrtati iznad broja 1 i spojiti ga linijom sa 1, a ne\'cemo
  ga spojiti linijom ni sa jednim drugim elementom na slici.
\end{minipage}
\hfill
\begin{minipage}{2.5in}
  \begin{center}
    \input ft07-r6e.pgf
  \end{center}
\end{minipage}

\bigskip

\noindent
\begin{minipage}{3in}
  Broj 6 je deljiv brojevima 2 i 3. Zato \'cemo broj 6 nacrtati iznad brojeva 2 i 3 i spojiti ga linijom i sa 2 i sa 3.
  Broj 6 je deljiv i brojem 1, ali smo se dogovorili da ne crtamo linije ukoliko postoji put koji spaja uporedive elemente i kre\'ce
  se samo navi\v se.
\end{minipage}
\hfill
\begin{minipage}{2.5in}
  \begin{center}
    \input ft07-r6f.pgf
  \end{center}
\end{minipage}

\bigskip

\noindent
\begin{minipage}{3in}
  Od svih navedenih brojeva broj 7 je deljiv samo brojem 1. Zato \'cemo broj 7 nacrtati iznad broja 1 i spojiti ga linijom sa 1, a ne\'cemo
  ga spojiti linijom ni sa jednim drugim elementom na slici.
\end{minipage}
\hfill
\begin{minipage}{2.5in}
  \begin{center}
    \input ft07-r6g.pgf
  \end{center}
\end{minipage}

\bigskip

\noindent
\begin{minipage}{3in}
  Od svih navedenih brojeva broj 8 je deljiv samo brojevima 1, 2 i 4. Zato \'cemo broj 8 nacrtati iznad broja 4 i spojiti ga linijom
  samo sa brojem 4. To je, prema dogovoru, dovoljno da se nazna\v ci da je 8 deljiv i sa 4, i sa 2, i sa 1.
\end{minipage}
\hfill
\begin{minipage}{2.5in}
  \begin{center}
    \input ft07-r6h.pgf
  \end{center}
\end{minipage}

\bigskip

\noindent
\begin{minipage}{3in}
  Sli\v cno tome, broj 9 je deljiv brojevima 1 i 3, pa \'cemo broj 9 nacrtati iznad broja 3 i spojiti ga linijom
  samo sa brojem 3.
\end{minipage}
\hfill
\begin{minipage}{2.5in}
  \begin{center}
    \input ft07-r6i.pgf
  \end{center}
\end{minipage}

\bigskip

\noindent
\begin{minipage}{3in}
  Kona\v cno, broj 10 je deljiv brojevima 2 i 5, pa \'cemo ga nacrtati iznad brojeva 2 i 5, i spojiti ga linijom
  samo sa oba. Time je Haseov dijagram u ovom primeru kompletiran.
\end{minipage}
\hfill
\begin{minipage}{2.5in}
  \begin{center}
    \input ft07-r6z.pgf
  \end{center}
\end{minipage}

\bigskip

\noindent
\begin{minipage}{3in}
  \begin{EX}\label{ft07.ex.P123}
    Pored je dat Haseov dijagram relacije poretka $\subseteq$ na skupu
    $\calP(\{a, b, c\})$.
  \end{EX}
\end{minipage}
\hfill
\begin{minipage}{2.5in}
  \begin{center}
    \input ft07-r7.pgf
  \end{center}
\end{minipage}

\bigskip

\noindent
\begin{minipage}{3in}
  \begin{EX}
    Pored je dat Haseov dijagram relacije poretka $\le$ na skupu
    $\{1, 2, 3, 4, 5\}$. Svaki linearni poredak ima ovakav Haseov dijagram.
    Odatle ime -- \emph{linearni poredak}.
  \end{EX}
\end{minipage}
\hfill
\begin{minipage}{2.5in}
  \begin{center}
    \input ft07-r8.pgf
  \end{center}
\end{minipage}

\bigskip

\emph{Najmanji element} relacije poretka na skupu $A$ je element skupa $A$ koji je uporediv sa svim elementima skupa $A$
i manji od svih njih. Sli\v cno, \emph{najve\'ci element} relacije poretka na skupu $A$ je element skupa $A$
koji je uporediv sa svim elementima skupa $A$ i ve\'ci od svih njih. Preciznije, neka je $\sqsubseteq$ relacija
poretka na skupu $A$. Element $a \in A$ je \emph{najmanji} element ako
$$
  (\forall x \in A) \; a \sqsubseteq x,
$$
a element $b \in A$ je \emph{najve\'ci} element ako
$$
  (\forall x \in A) \; x \sqsubseteq b.
$$

\begin{EX}\label{ft07.ex.delj-2-12}
  Uobi\v cajena relacija poretka $\le$ na skupu $\ZZ$ nema ni najmanji ni najve\'ci element.

  Relacija poretka koja je prikazana u Primeru~\ref{ft07.ex.P123} ima i najmanji i najve\'ci element.

  Relacija poretka koja je prikazana u Primeru~\ref{ft07.ex.deljivost-1-10} ima najmanji element, ali nema najve\'ci.
\end{EX}

\bigskip

\noindent
\begin{minipage}{3in}
  \hspace*{1.5em}Relacija poretka $|$ (deljivost) na skupu $\{2, 3, 4, 6, 12\}$ je prikazana pored. Ona ima najve\'ci element, ali
  nema najmanji element.
\end{minipage}
\hfill
\begin{minipage}{2.5in}
  \begin{center}
    \input ft07-r10.pgf
  \end{center}
\end{minipage}

\bigskip

Pogledajmo jo\v s jednom relacije poretka iz Primera~\ref{ft07.ex.deljivost-1-10} i~\ref{ft07.ex.delj-2-12}:
\begin{center}
    \input ft07-r6z.pgf
    \hspace*{2cm}
    \input ft07-r10.pgf
\end{center}
Vidimo da prva relacija ima najmanji element, ali nema najve\'ci element, dok druga relacija ima najve\'ci element, ali nema
najmanji element. Ako pogledamo element 10 u prvoj relaciji, vidimo da je on nekako ``pri vrhu'' i da nije najve\'ci element samo
zato \v sto nije uporediv sa elementima kao \v sto su 6, 7, 8, 9. Za takav element relacije poretka ka\v zemo da je
\emph{maksimalan element}. Sli\v cno, element 2 druge relacije poretka je nekako ``pri dnu'' i to nije najmanji element
samo zato \v sto nije uporediv sa~3. Za takav element relacije poretka ka\v zemo da je \emph{minimalan element}.

Preciznije, neka je $\sqsubseteq$ relacija poretka na skupu $A$. Element $a \in A$ je \emph{maksimalan element} ako
$$
  \lnot(\exists x \in A)(a \ne x \land a \sqsubseteq x),
$$
odnosno, ako iznad njega nema drugih elemenata. Sli\v cno, element $b \in A$ je \emph{minimalan element} ako
$$
  \lnot(\exists x \in A)(x \ne b \land x \sqsubseteq b),
$$
odnosno, ako ispod njega nema drugih elemenata. Relacija poretka mo\v ze, a ne mora imati maksimalne i minimalne elemente,
minimalnih i maksimalnih elemenata mo\v ze biti vi\v se, a mo\v ze se desiti da je neki element istovremeno i maksimalan
i minimalan element relacije pretka, kao \v sto pokazuje slede\'ci primer.

\bigskip

\noindent
\begin{minipage}{3in}
  \begin{EX}
    Relacija poretka $|$ (deljivost) na skupu $\{2, 3, 4, 5, 6\}$ je prikazana pored. Minimalni elementi ove relacije
    su 2, 3 i 5 jer ispod njih nema drugih elemenata. Maksimalni elementi ove relacije su 4, 6 i 5 jer iznad njih nema
    drugih elemenata. Vidimo da je 5 istovremeno i minimalni i maksimalni element.
  \end{EX}
\end{minipage}
\hfill
\begin{minipage}{2.5in}
  \begin{center}
    \input ft07-r11.pgf
  \end{center}
\end{minipage}

\begin{THM}
  Neka je $\sqsubseteq$ relacija poretka na skupu $A$.
  \begin{enumerate}
  \item
    Ako relacija poretka $\sqsubseteq$ ima najve\'ci element, onda je to ujedno i njen jedinstveni maksimalni element.
  \item
    Ako relacija poretka $\sqsubseteq$ ima najmanji element, onda je to ujedno i njen jedinstveni minimalni element.
  \end{enumerate}
\end{THM}


\section{Apstrakcija relacije poretka u programiranju}

Pretpostavimo da treba da napravimo jednostavan model u kome o studentima pamtimo slede\'ce podatke:
{\small
\begin{verbatim}
    String ime, prezime;
    int brIdx; // broj indeksa
    double prosek;
\end{verbatim}
}
Pretpostavimo da imamo i dodatni zahtev da sistem treba da bude organizovan tako da se liste studenata mogu
lako sortirati. Programski jezik Java ima
klasu \verb`Collections` koja implementira metod \verb`sort`, ali programer mora da objasni sistemu
po kom kriterijumu lista treba da bude sortirana. Drugim re\v cima, programer mora da opi\v se relaciju poretka
koju \v zeli da primeni na instance klase \verb`Student` kako bi metod \verb`Collections.sort` mogao
da sortira elemente.

Apstrakcija relacije linearnog ure\dj enja u Javi je realizovana kao parametrizovani interfejs \verb`Comparable`, gde
parametar interfejsa predstavlja tip objekata koje \'cemo porediti. Ovaj interfejs sadr\v zi metod
{\small
\begin{verbatim}
    int compareTo(T obj)
\end{verbatim}
}
\noindent
koji treba da vrati $-1$ ukoliko je \verb`this` strogo manje od \verb`obj`, treba da vrati $0$ ukoliko su \verb`this` i \verb`obj`
jednaki, i treba da vrati $1$ ukoliko je \verb`this` strogo ve\'ce od \verb`obj`. U na\v sem modelu \'ce, zato, klasa
\verb`Student` implementirati interfejs
{\small
\begin{verbatim}
    Comparable<Student>
\end{verbatim}
}
\noindent
(sa parametrom \verb`Student` jer \v zelimo
da poredimo studente) i implementira\'ce metod \verb`compareTo` tako da studente poredi prvo po prezimenu, potom
po imenu ako su prezimena ista, i kona\v cno po broju indeksa ako su imena i prezimena ista. Klasa \verb`Student`, zato,
izgleda ovako:
{\small
\begin{verbatim}
  public class Student implements Comparable<Student> {
    String ime, prezime;
    int brIdx; // broj indeksa
    double prosek;
    Student(String ime, String prezime, int brIdx, double prosek){
      this.ime = ime; this.prezime = prezime;
      this.brIdx = brIdx; this.prosek = prosek;
    }
    
    @Override
    public int compareTo(Student s){
      // poredimo prezimena studenata
      int k = prezime.compareTo(s.prezime);
      if(k != 0) return k;
      // studenti imaju isto prezime, pa poredimo imena
      k = ime.compareTo(s.ime);
      if(k != 0) return k;
      // studenti imaju isto ime i prezime,
      // pa poredimo brojeve indeksa
      if(brIdx < s.brIdx) return -1; 
      if(brIdx > s.brIdx) return  1;
      return 0;
    }
  }
\end{verbatim}
}

\noindent
Ovako struktuiran model nam omogu\'cuje da automatski sortiramo liste studenata, kako pokazuje slede\'ci primer:
{\small
\begin{verbatim}
  import java.util.ArrayList;
  import java.util.Collections;
  public class StudentSort{
    public static void main(String args[]){
      ArrayList<Student> lista = new ArrayList<Student>();
      lista.add(new Student("Mirjana","Mirkovic", 34218, 7.25));
      lista.add(new Student("Petar","Petrovic",  12324, 8.45));
      lista.add(new Student("Jovan","Jovanovic", 23256, 9.05));
      lista.add(new Student("Mirko","Mirkovic", 34213, 7.25));
      lista.add(new Student("Jovanka","Jovanovic", 23257, 9.27));
        
      Collections.sort(lista);
      for(Student s : lista){
        System.out.println(s.brIdx + " " + s.prezime + ", " + s.ime);
      }
    }
  }
\end{verbatim}
}




